\chapter{Temporal Reconstruction of the Velocity Field}\label{ch:temporal}

The integration schemes of \cref{ch:integration} require the velocity not only at
arbitrary positions but at arbitrary times --- at the half-step and full-step instants
of the Runge--Kutta stages --- whereas the model stores it only at discrete snapshots.
The engine bridges this gap by reconstructing the velocity in time between snapshots.
This chapter describes the two modes of operation (one snapshot per step, or
sub-stepping with interpolation), the linear and monotone-cubic reconstructions, and the
central consequence that the temporal reconstruction can cap the order of the spatial
scheme. The handling of the snapshot times themselves --- including the shift applied to
time-averaged forcing --- belongs to the forcing reader and is described in
\cref{ch:input}.

\section{Discrete and interpolated modes}\label{sec:temp-modes}

The key \conf{temporal\_interp} selects between two modes. In the \emph{discrete} mode (the default),
each integration step advances by one stored snapshot; the displacement uses the step \conf{numerics.dt},
which must therefore equal the record spacing (taken from the file's time axis when it is not set; a
different step, or irregular record times such as calendar months, are refused when the run starts), and
the Runge--Kutta stages are
still ordered in time, the next snapshot being supplied so that the stages see the correct time
progression across the step. In the \emph{interpolated} mode, the time step may be smaller than the
snapshot interval, and the velocity at each stage time is reconstructed from the bracketing snapshots.
Sub-stepping in the interpolated mode is the means by which a fine step can be taken against coarsely
sampled forcing, for example a sub-daily step against daily-mean fields. The interpolated mode needs a
physical time axis: a file without a decodable time coordinate is refused unless the snapshot spacing
is given as \conf{numerics.file\_dt}.

\section{Linear reconstruction}\label{sec:temp-linear}

The default reconstruction is linear in time. For a time $t$ falling between snapshots at $t_a$ and $t_b$
carrying velocity fields $\velvec_a$ and $\velvec_b$, the reconstructed field is
\begin{equation}\label{eq:temp-linear}
  \velvec(t) = (1-\interpfrac)\,\velvec_a + \interpfrac\,\velvec_b, \qquad
  \interpfrac = \frac{t-t_a}{t_b-t_a}.
\end{equation}
A step starting at $t$ between two snapshots receives the field at its start, with
$\interpfrac_0=(t-t_a)/(t_b-t_a)$, and a \emph{target} snapshot $\velvec_T$ toward which its stages
move: the next snapshot, $t_b$, when integrating forward, and the previous one, $t_a$, when integrating
backward; a backward step that starts exactly on a snapshot brackets downward, so it has a real target
rather than none. The stage with time fraction $c_k$ (\cref{sec:int-schemes}) samples
\begin{equation}\label{eq:temp-stage}
  \velvec_k=(1-\interpfrac_k)\,\velvec(t)+\interpfrac_k\,\velvec_T,\qquad
  \interpfrac_k=\frac{c_k\,\dt}{\Delta t_T},
\end{equation}
with $\Delta t_T$ the time from the step's start to the target snapshot; \cref{eq:temp-stage} is
\cref{eq:temp-linear} evaluated at $t\pm c_k\dt$. The engine never forms $\velvec(t)$ as a whole field
unless a consumer asks for it: the velocity handed to a step between two snapshots is a lazy view that
blends a field on its first use, and the fused kernels blend each element or vertex they read
(\cref{eq:fused-nested}).

A step longer than the time to its target snapshot, $\dt>\Delta t_T$, would place its later stages
beyond the bracketing pair. Such a step runs on a steady velocity --- the field at its start --- and the
engine warns once per realization and reports the number of such steps at the end of a verbose run.
This happens at the end of the record, on a start between snapshots, and whenever $\dt$ does not
divide the snapshot interval: every step that straddles a snapshot then runs steady.

\section{The order cap, and monotone-cubic reconstruction}\label{sec:temp-pchip}

A linear reconstruction in time is itself only second-order accurate, and this caps the
\emph{time}-truncation error of the whole integration at order $\dt^2$ irrespective of
the spatial scheme. The practical consequence is important: under linear reconstruction,
RK4 collapses onto RK2 --- the fourth-order spatial scheme delivers no benefit --- unless
the step is taken much smaller than the snapshot interval. The remedy, short of finer
sub-stepping, is a higher-order reconstruction in time, selected by
\conf{temporal\_reconstruction}${}={}$\conf{pchip}.

The PCHIP reconstruction (\cref{fig:temporal}) fits a monotone piecewise-cubic Hermite
interpolant \citep{fritsch1980monotone} through four consecutive snapshots
$t_{n-1},t_n,t_{n+1},t_{n+2}$, evaluating it within the central interval
$[t_n,t_{n+1}]$. On that interval, with local coordinate $s=(t-t_n)/h_1$ and
$h_1=t_{n+1}-t_n$, the value is the Hermite combination
\begin{equation}\label{eq:pchip}
  f(t) = h_{00}(s)\,f_n + h_{10}(s)\,h_1 m_n + h_{01}(s)\,f_{n+1} + h_{11}(s)\,h_1 m_{n+1},
\end{equation}
with the standard cubic basis $h_{00}=2s^3-3s^2+1$, $h_{10}=s^3-2s^2+s$,
$h_{01}=-2s^3+3s^2$, $h_{11}=s^3-s^2$. The slope at an interior snapshot is the weighted harmonic
mean of the two adjacent secants $d_-$ and $d_+$ over intervals $h_-$ and $h_+$
\citep{fritsch1984method},
\begin{equation}\label{eq:pchip-slope}
  m=\frac{w_1+w_2}{w_1/d_-+w_2/d_+},\qquad w_1=2h_++h_-,\quad w_2=h_++2h_-,
\end{equation}
and zero wherever the two secants differ in sign or either vanishes. This makes the interpolant
monotone between samples and prevents the spurious overshoots that an unconstrained cubic would
introduce. The property holds for each component separately: the cubic does not preserve linear
relations between components or between snapshots --- a vanishing divergence, rotation equivariance, a
window mean --- which the linear reconstruction does preserve. Because the reconstruction is now cubic in time, the fourth-order spatial
scheme is no longer masked: RK4 separates from RK2 at coarse snapshot cadence. The PCHIP
stencil requires four consecutive snapshots, so on the first and last interval of the
record, and across a cyclic seam, the engine falls back to the linear reconstruction, with one warning
per realization. The reconstruction applies to the horizontal and vertical velocity and the sea-surface
height; the diffusivities stay linear in time, and sampled scalars are re-interpolated linearly to
the record time (\cref{sec:samp-what}). PCHIP is active
only with \conf{temporal\_interp: true} (otherwise it warns that it has no effect) and only for the
multistage schemes RK2 and RK4; the Euler step has no stage times to reconstruct.

The two-dimensional least-squares kernels and the three-dimensional barycentric RK2 kernel evaluate the
cubic inline for each element or vertex they read (\cref{ch:fused}). Every other configuration
reconstructs whole fields at each distinct stage time --- the two midpoint stages of RK4 share one ---
with a parallel loop over the field that is bit-identical to the plain array expression; one such blend
of a global field takes about \SI{46}{ms} at eight threads.

\begin{figure}[t]
\centering
\begin{tikzpicture}[scale=1.0]
  % axes
  \draw[-{Stealth[length=5pt]},gray] (-0.2,0)--(7.2,0) node[right,font=\footnotesize,black]{time};
  \draw[-{Stealth[length=5pt]},gray] (0,-0.2)--(0,3.0) node[above,font=\footnotesize,black]{$\velvec$};
  % snapshot points
  \coordinate (s0) at (0.6,1.0);
  \coordinate (s1) at (2.6,2.3);
  \coordinate (s2) at (4.6,1.5);
  \coordinate (s3) at (6.6,2.5);
  % PCHIP smooth curve through all four
  \draw[laggreen,line width=1.1pt,smooth] plot coordinates {(0.6,1.0) (1.6,1.75) (2.6,2.3) (3.6,2.0) (4.6,1.5) (5.6,1.8) (6.6,2.5)};
  % linear chord on central interval [t_n, t_{n+1}] = s1..s2
  \draw[lagblue,line width=1.0pt,densely dashed] (s1)--(s2);
  % central interval shading
  \draw[gray!55,densely dotted] (2.6,0)--(2.6,2.3);
  \draw[gray!55,densely dotted] (4.6,0)--(4.6,1.5);
  % query substep
  \coordinate (q) at (3.6,0);
  \draw[gray!55,densely dotted] (3.6,0)--(3.6,2.02);
  \node[font=\footnotesize,below] at (3.6,0) {$t$};
  \node[circle,fill=lagblue,inner sep=1.3pt] at (3.6,1.9) {};
  \node[circle,fill=laggreen,inner sep=1.3pt] at (3.6,2.0) {};
  % snapshot markers + labels
  \foreach \p in {s0,s1,s2,s3} \node[interfacept] at (\p) {};
  \node[font=\footnotesize,below] at (0.6,-0.02) {$t_{n-1}$};
  \node[font=\footnotesize,below] at (2.6,-0.02) {$t_{n}$};
  \node[font=\footnotesize,below] at (4.6,-0.02) {$t_{n+1}$};
  \node[font=\footnotesize,below] at (6.6,-0.02) {$t_{n+2}$};
  \node[font=\footnotesize,laggreen!60!black] at (5.7,2.25) {PCHIP (4-point)};
  \node[font=\footnotesize,lagblue!70!black] at (3.05,2.45) {linear (2-point)};
\end{tikzpicture}
\caption{Temporal reconstruction within the interval $[t_n,t_{n+1}]$. The linear
reconstruction (dashed) joins the two bracketing snapshots; the monotone-cubic PCHIP
reconstruction (solid) uses four consecutive snapshots and follows the curvature of the
sampled series. At a sub-step time $t$ the two differ, and only the cubic lets a
fourth-order spatial scheme realize its order. PCHIP is monotone, so it introduces no
overshoot between samples.}
\label{fig:temporal}
\end{figure}

\section{Read-ahead}\label{sec:temp-prefetch}

Reconstruction in time requires several snapshots to be in memory at once --- two for the linear scheme,
four for PCHIP. Snapshots stay in a small cache as long as the current bracket or stencil needs them, so
a PCHIP run reads one new snapshot each time its bracket advances. The key \conf{prefetch\_size} controls
how many consecutive snapshots are read in a single bulk operation: a value above one reads a window ahead
in the direction of travel, reducing the number of separate reads on large meshes where input dominates
the cost (\cref{ch:performance}), while a value of one reads one snapshot at a time; the trajectories are
bit-identical either way. Read-ahead is disabled under cyclic forcing, where the wrap-around would
otherwise thrash the cache, and when snapshots are shared between processes through
\conf{shared\_fields\_dir} (\cref{ch:input}).

\codeanchor{The reconstructions are in \srcfile{numerics/temporal\_reconstruct.py} (the slopes
\fn{\_fc\_slope}, the Hermite evaluation and the parallel whole-field blend \fn{\_pchip\_blend4}); the
lazy linear view \fn{LazyLerpField} and its blend are in \srcfile{core/io.py}; the stage fields of the
unfused paths are built by \fn{\_make\_stage\_vels} and \fn{make\_stage\_vels\_recon}; the bracket and
stencil resolution (\fn{\_resolve\_time\_bracket}, \fn{\_resolve\_recon\_stencil}), the step payload
(\fn{\_build\_vel\_payload}) and the read-ahead cache (\fn{\_VelCache}) are in \srcfile{tracking.py}.}
